\section{Introduction}

Reinforcement learning (RL) depends critically on the reward function used to guide policy optimization. A well-designed reward can accelerate learning and align the policy with the intended task, whereas an incomplete or misspecified reward may induce sparse exploration, unstable training, reward hacking, or behaviors that maximize the designed signal while failing the real task objective. Designing such rewards typically requires expert knowledge of the environment dynamics, the task semantics, and the failure modes of the learner.

Large language models (LLMs) provide a promising alternative for reward design because they can interpret task descriptions, reason over environment code, and synthesize executable programs. Recent work such as Eureka demonstrates that coding LLMs can generate reward functions that support reinforcement learning across diverse control tasks \cite{ma2023eureka}. Other methods, including Text2Reward and language-model-based reward design, further show that natural language can be used to specify desired behavior or shape dense rewards \cite{xie2023text2reward,kwon2023rewarddesign}. These works suggest a new direction: using LLMs not merely as policy planners, but as reward designers.

However, LLM-based reward generation still faces three limitations. First, a reward generated in a single pass is often insufficient. Reward functions must be evaluated through training and then revised when they produce poor exploration, unstable learning, or undesired shortcuts. Second, feedback used for reward improvement is often entangled with official task rewards, human correction, or manually designed success metrics. Such feedback can be effective, but it weakens the claim that the reward is discovered independently of the task's oracle reward. Third, direct Python reward generation is flexible but difficult to validate, edit, attribute, and reuse across iterations.

This paper proposes \methodname{}: a memory-reflective reward schema adaptation framework for LLM-driven reward self-evolution without oracle reward feedback. Instead of asking an LLM to repeatedly output unrestricted Python reward code, \methodname{} represents a reward function as a structured reward schema. A schema consists of formula-based reward components, event predicates, weights, clips, and semantic roles. This representation preserves expressiveness while enabling canonicalization, validation, component attribution, and localized edits.

The feedback mechanism in \methodname{} is deliberately non-oracle. During policy training, the environment reward is replaced by the schema reward. The official reward, if available, is not used to guide reward search. After training, \methodname{} builds feedback from rollout trajectories, reward component attribution, trajectory summaries, and behavior-level diagnostics. This feedback tells the system whether the current reward is dominated by a single component, whether the policy exploits sparse events, whether progress is insufficient, or whether the observed behavior mismatches the task description.

A key design of \methodname{} is memory-reflective iteration. The system maintains three levels of memory: raw memory cards that record exact edit and outcome traces, distilled lesson cards that summarize reusable experience, and outcome lessons that compare before/after effects of committed edits. A ReflectionAgent examines these memory layers and decides which evidence to reuse, ignore, or treat as conflicting. An EditAgent then turns the reflection strategy into a concrete schema-level edit plan. Finally, a SchemaEngine validates the edit plan, applies the committed edits transactionally, canonicalizes the next schema, and passes it to the next training iteration.

The main contributions of this work are:
\begin{itemize}
    \item We introduce a non-oracle feedback mechanism for LLM-driven reward search. The method improves rewards using rollout behavior, component attribution, and task-proxy diagnostics rather than official environment rewards.
    \item We propose a memory-reflective reward iteration mechanism. The ReflectionAgent selectively reasons over raw memory, distilled lessons, and measured outcome lessons before reward editing.
    \item We formulate reward generation as schema-level self-evolution. The reward schema acts as a safe, editable intermediate representation that supports canonicalization, validation, attribution, and iterative modification.
\end{itemize}
